% GENERATED FILE. Edit canonical lessons, then run npm run generate:books.
% canonical-source-hash: fnv1a64:197ce29d53d33c30
% canonical-lessons: ZH-C151-kandao, ZH-W151-dao-arrive, ZH-C151-result, ZH-C151-xuehui, ZH-C151-tingdong, ZH-C151-kandong, ZH-C151-zuowan, ZH-C151-practice

\chapter{Seen and Understood}
\label{ch:zh-seen-and-understood}
\hlchaptermodality{151}

\begin{chapteropening}
\textbf{By the end of this chapter:} \emph{I can say that I saw, heard, understood or finished something.}

A verb and its result: seeing, learning how, understanding by ear and by eye, and finishing, each negated with \zh{没}, built on one new written character: \zh{到}.
\end{chapteropening}

\section[\texorpdfstring{\zh{看到} (kàndào)}{看到 (kàndào)}]{\zh{看到} (kàndào) — to see; to catch sight of}

\label{lesson:ZH-C151-kandao}



\begin{quote}
A new chapter, and it opens on a verb that reaches what it aims at.

\begin{itemize}
\raggedright
  \item \emph{Recall:} say \zh{你去过吗} — \emph{nǐ qù guo ma}, and answer it — \textbf{R1}, the chapter you have finished
  \item \emph{Recall:} say \emph{cónglái}, ever — \textbf{R1}
  \item \emph{Recall:} say \zh{看见} — \emph{kànjiàn}, to catch sight of — \textbf{R4}, recognition at distance
  \item \emph{Recall:} say \emph{dào}, to arrive — \textbf{R3}, durable
\end{itemize}
\end{quote}

\subsection*{You'll want to know: \zh{看到}}
\textbf{\zh{看到}} — \emph{kàndào} — \textquotedblleft{}to see; to catch sight of\textquotedblright{}.

\begin{quote}
\textbf{\zh{我看到他了}} \emph{wǒ kàndào tā le} — I saw him.
\end{quote}

\begin{quote}
\textbf{\zh{我没看到}} \emph{wǒ méi kàndào} — I didn't spot it.
\end{quote}

The next lesson puts the second character in your hand. It is the \emph{dào} you already say for arriving.

\begin{cousinweb}
\textbf{\zh{看到}} (\emph{kàndào}) is built from \zh{看} 'to look' and \zh{到} 'to arrive; to reach'. \zh{看} is inherited from Middle Chinese \emph{khanH} (Guangyun, in Baxter's transcription), which the Guangyun lists beside a level-tone reading glossed 'to look'; \zh{到} is inherited from \emph{tawH}, glossed 'to arrive'.
\end{cousinweb}

\begin{grammarlens}[title={What you've built}]
\zh{看} is the looking; \zh{到} says the looking reached something.

\zh{看见}, which you learned long ago as one word, is made the same way: \zh{看}, then \zh{见}.
\end{grammarlens}

\subsection*{Your turn}
\begin{itemize}
\raggedright
  \item \emph{Say it:} \emph{kàndào}
  \item \emph{Say it:} \emph{kàndào}, once more
  \item \emph{Say it:} \zh{我看到他了} — \emph{wǒ kàndào tā le} — I saw him.
  \item \emph{Say it:} \zh{我没看到} — \emph{wǒ méi kàndào} — I didn't spot it.
  \item \emph{Recall:} say \zh{看见}, then \zh{看到}
\end{itemize}

\begin{tcolorbox}[breakable,colback=teal!4,colframe=teal!35!black,title={Before you move on}]
What does \zh{看到} mean? (\textquotedblleft{}To see; to catch sight of\textquotedblright{}.) Which word you know is made the same way? (\textbf{\zh{看见}}.)
\end{tcolorbox}
\section[\texorpdfstring{\zh{到} (dào)}{到 (dào)}]{\zh{到} — \zh{至} on the left, then the blade of \zh{前}}

\label{lesson:ZH-W151-dao-arrive}



\begin{quote}
The word first, then the character inside it.

\begin{itemize}
\raggedright
  \item \emph{Recall:} say \emph{kàndào} — \textbf{R1}, one lesson back, before it decays
\end{itemize}

\begin{quote}
\zh{在} \zh{前}
\end{quote}

\emph{Read it:} each one aloud, with nothing said first — \emph{zài}, \emph{qián}
\end{quote}

\subsection*{Script — \zh{至} on the left, then the blade of \zh{前}}
\hlblockfigure{\detokenize{figures/ZH-W151-dao-arrive-filmstrip.pdf}}{How \zh{到} is written, stroke by stroke}

\begin{quote}
\zh{到}
\end{quote}

\textbf{Two pieces, and both have parts you write.}

\textbf{\zh{至}, six strokes}, on the left: a bar across the top; a stroke that falls to the left and turns right; a dot at its end; then \zh{土}, as at the foot of \zh{在}: a short bar, a vertical through it, and the bottom bar, rising a little to the right.

\textbf{\zh{刂}, two strokes}, on the right, as on the right of \zh{前}: a short vertical, then a long vertical that hooks to the left at its foot.

6 + 2 = eight.

\subsection*{Your turn}
\begin{itemize}
\raggedright
  \item \emph{Write it:} \zh{到}, all eight, with the model in front of you
  \item \emph{Cover:} write it again from memory, and count the pieces, not the strokes
\end{itemize}

\begin{tcolorbox}[breakable,colback=teal!4,colframe=teal!35!black,title={Before you move on}]
How many strokes in \zh{到}? (\textbf{Eight}.) Which character you write ends on the same \zh{土}? (\textbf{\zh{在}}.) And on the same two right-hand strokes? (\textbf{\zh{前}}.)

Source: \href{https://raw.githubusercontent.com/chanind/hanzi-writer-data/68d10a4b21150cae5e1ebbd223eed289cf32d90c/data/\%E5\%88\%B0.json}{Hanzi Writer Data \zh{到}, pinned stroke paths and medians}.
\end{tcolorbox}
\section[\texorpdfstring{\zh{看到} — \zh{没看到} (kàndào — méi kàndào)}{看到 — 没看到 (kàndào — méi kàndào)}]{\zh{看到} — \zh{没看到}: a verb and its result}

\label{lesson:ZH-C151-result}



\begin{quote}
Two verbs in a row, and the second one says how it turned out.

\begin{itemize}
\raggedright
  \item \emph{Read it:} \zh{看到}, aloud — \emph{kàndào}
  \item \emph{Recall:} say \zh{看见} — \emph{kànjiàn}
  \item \emph{Recall:} say \zh{我没去} — \emph{wǒ méi qù}, I didn't go
\end{itemize}
\end{quote}

\begin{grammarlens}[title={Grammar Lens: a verb, then what came of it}]
\begin{quote}
\textbf{A second word straight after a verb can say how it turned out. \zh{看} is looking; \zh{看到} is looking that reached something.}
\end{quote}

\begin{quote}
\textbf{\zh{我看了} \zh{没看到}} \emph{wǒ kàn le, méi kàndào} — I looked, but I didn't spot it.
\end{quote}

The negative is \zh{没}, because the result did not come about:

\begin{quote}
\textbf{\zh{我没看到他}} \emph{wǒ méi kàndào tā} — I didn't see him. \textbf{\zh{你看到了吗}} \emph{nǐ kàndào le ma} — Did you see it?
\end{quote}

\zh{看见} is made the same way, \zh{看} and its result \zh{见}. The next four lessons add four more pairs like these, each a verb and how it ended.
\end{grammarlens}

\subsection*{Your turn}
Say these aloud:

\begin{itemize}
\raggedright
  \item \zh{我看了} \zh{没看到} — \emph{wǒ kàn le, méi kàndào} — I looked, but I didn't spot it.
  \item \zh{你看到了吗} — \emph{nǐ kàndào le ma} — Did you see it?
  \item \zh{看到了} — \emph{kàndào le} — Yes, I saw it.
  \item \zh{我没看到他} — \emph{wǒ méi kàndào tā} — I didn't see him.
\end{itemize}

\begin{tcolorbox}[breakable,colback=teal!4,colframe=teal!35!black,title={Before you move on}]
What does the second word of \zh{看到} say? (\textbf{How the looking turned out}.) Which negative goes in front? (\textbf{\zh{没}: \zh{没看到}}.) Which old word is made the same way? (\textbf{\zh{看见}}.)
\end{tcolorbox}
\section[\texorpdfstring{\zh{学会} (xuéhuì)}{学会 (xuéhuì)}]{\zh{学会} (xuéhuì) — to learn, and now be able to}

\label{lesson:ZH-C151-xuehui}



\begin{quote}
Before the new one: read the last lesson's pair, then \zh{会} as a skill.

\begin{itemize}
\raggedright
  \item \emph{Read it:} \zh{看到} — \zh{没看到}, aloud — \emph{kàndào}, \emph{méi kàndào}
  \item \emph{Recall:} say \zh{我会开车} — \emph{wǒ huì kāichē}, I can drive — \textbf{R3}, durable
  \item \emph{Recall:} say \emph{xué}, to study — \textbf{R4}, recognition at distance
\end{itemize}
\end{quote}

\subsection*{You'll want to know: \zh{学会}}
\textbf{\zh{学会}} — \emph{xuéhuì} — \textquotedblleft{}to learn (and now be able to)\textquotedblright{}.

\begin{quote}
\textbf{\zh{我学会开车了}} \emph{wǒ xuéhuì kāichē le} — I've learned to drive.
\end{quote}

\begin{cousinweb}
\textbf{\zh{学会}} (\emph{xuéhuì}) is built from \zh{学} 'to study' and \zh{会} 'can'. Both are inherited from Middle Chinese, \emph{haewk} and \emph{hwajH} (Guangyun, in Baxter's transcription), glossed 'to awaken to; to learn' and 'to come together'. \zh{会} as 'can' is a later sense than that gloss, assumed here, not quoted.
\end{cousinweb}

\begin{grammarlens}[title={What you've built}]
\zh{学} is the studying; \zh{会} is the result: now you can. With \zh{了} at the end it is news: \zh{我学会开车了}.

Without the result, \textbf{\zh{我学开车}} \emph{wǒ xué kāichē} says only that you are studying it.
\end{grammarlens}

\subsection*{Your turn}
\begin{itemize}
\raggedright
  \item \emph{Say it:} \emph{xuéhuì}
  \item \emph{Say it:} \emph{xuéhuì}, once more
  \item \emph{Say it:} \zh{我学会开车了} — \emph{wǒ xuéhuì kāichē le} — I've learned to drive.
  \item \emph{Say it:} \zh{我没学会} — \emph{wǒ méi xuéhuì} — I haven't learned how.
  \item \emph{Recall:} say \zh{我会开车}, then \zh{我学会开车了}
\end{itemize}

\begin{tcolorbox}[breakable,colback=teal!4,colframe=teal!35!black,title={Before you move on}]
What does \zh{学会} mean? (\textquotedblleft{}To learn, and now be able to\textquotedblright{}.) Which part is the result? (\textbf{\zh{会}}.)
\end{tcolorbox}
\section[tīngdǒng]{tīngdǒng — to understand what you hear}

\label{lesson:ZH-C151-tingdong}



\begin{quote}
Before the new one: last lesson's word, then the two verbs this one is made of.

\begin{itemize}
\raggedright
  \item \emph{Recall:} say \zh{学会} — \emph{xuéhuì} — \textbf{R1}, one lesson back
  \item \emph{Recall:} say \emph{tīng}, to listen — \textbf{R3}, durable
  \item \emph{Recall:} say \emph{dǒng}, to understand — \textbf{R3}
\end{itemize}
\end{quote}

\subsection*{You'll want to know: tīngdǒng}
\emph{tīngdǒng} — \textquotedblleft{}to understand (on hearing)\textquotedblright{}.

\begin{quote}
\emph{wǒ tīngdǒng le} — I understood what was said.
\end{quote}

\begin{quote}
\emph{nǐ tīngdǒng le ma} — Did you catch that?
\end{quote}

Its characters come later; for now it is a word for the ear.

\begin{cousinweb}
\emph{tīngdǒng} is built from \emph{tīng} 'to listen' and \emph{dǒng} 'to understand', the two verbs you know. The character for \emph{tīng} is inherited from Middle Chinese \emph{theng} (Guangyun, in Baxter's transcription), glossed 'to hear'. The character for \emph{dǒng} has no reading in the Guangyun: no dictionary we use records an older form, so its history is not settled.
\end{cousinweb}

\begin{grammarlens}[title={What you've built}]
\emph{tīng} is the listening; \emph{dǒng} is the result. \zh{没} says it did not come: \emph{wǒ méi tīngdǒng}, I didn't catch it.

\emph{tīng} alone is only listening; \emph{tīngdǒng} is listening that worked.
\end{grammarlens}

\subsection*{Your turn}
\begin{itemize}
\raggedright
  \item \emph{Say it:} \emph{tīngdǒng}
  \item \emph{Say it:} \emph{tīngdǒng}, once more
  \item \emph{Say it:} \emph{nǐ tīngdǒng le ma} — Did you catch that?
  \item \emph{Say it:} \emph{wǒ méi tīngdǒng} — I didn't catch it.
  \item \emph{Recall:} say \emph{tīng}, then \emph{dǒng}, then \emph{tīngdǒng}
\end{itemize}

\begin{tcolorbox}[breakable,colback=teal!4,colframe=teal!35!black,title={Before you move on}]
What does \emph{tīngdǒng} mean? (\textquotedblleft{}To understand what you hear\textquotedblright{}.) Say you didn't catch it. (\emph{wǒ méi tīngdǒng}.)
\end{tcolorbox}
\section[kàndǒng]{kàndǒng — to understand what you read}

\label{lesson:ZH-C151-kandong}



\begin{quote}
Before the new one: last lesson's word, then reading a book.

\begin{itemize}
\raggedright
  \item \emph{Recall:} say \emph{tīngdǒng}, to understand what you hear — \textbf{R1}, one lesson back
  \item \emph{Recall:} say \zh{看书} — \emph{kàn shū}, to read — \textbf{R4}, recognition at distance
\end{itemize}
\end{quote}

\subsection*{You'll want to know: kàndǒng}
\emph{kàndǒng} — \textquotedblleft{}to understand (on reading or watching)\textquotedblright{}.

\begin{quote}
\emph{zhè běn shū wǒ kàndǒng le} — I understood this book.
\end{quote}

Its second character comes later; for now it is a word for the ear.

\begin{cousinweb}
\emph{kàndǒng} is built from \zh{看} 'to look; to read' and \emph{dǒng} 'to understand'. \zh{看} is inherited from Middle Chinese \emph{khanH} (Guangyun, in Baxter's transcription); the character for \emph{dǒng} has no reading in the Guangyun, so its history is not settled.
\end{cousinweb}

\begin{grammarlens}[title={What you've built}]
The thing can go first, as the topic: \emph{zhè běn shū}, then who, then the verb and its result.

\emph{tīngdǒng} for the ear, \emph{kàndǒng} for the eye: one result, two ways in.
\end{grammarlens}

\subsection*{Your turn}
\begin{itemize}
\raggedright
  \item \emph{Say it:} \emph{kàndǒng}
  \item \emph{Say it:} \emph{kàndǒng}, once more
  \item \emph{Say it:} \emph{zhè běn shū wǒ kàndǒng le} — I understood this book.
  \item \emph{Say it:} \emph{nǐ kàndǒng le ma} — Did you follow it?
  \item \emph{Recall:} say \emph{tīngdǒng}, then \emph{kàndǒng}
\end{itemize}

\begin{tcolorbox}[breakable,colback=teal!4,colframe=teal!35!black,title={Before you move on}]
What does \emph{kàndǒng} mean? (\textquotedblleft{}To understand what you read or watch\textquotedblright{}.) And by ear? (\emph{tīngdǒng}.)
\end{tcolorbox}
\section[zuòwán]{zuòwán — to finish (doing)}

\label{lesson:ZH-C151-zuowan}



\begin{quote}
Before the new one: last lesson's word, then the verb this one starts with.

\begin{itemize}
\raggedright
  \item \emph{Recall:} say \emph{kàndǒng}, to understand what you read — \textbf{R1}, one lesson back
  \item \emph{Recall:} say \zh{做饭} — \emph{zuò fàn}, to cook — \textbf{R3}, durable
\end{itemize}
\end{quote}

\subsection*{You'll want to know: zuòwán}
\emph{zuòwán} — \textquotedblleft{}to finish (doing)\textquotedblright{}.

\begin{quote}
\emph{wǒ zuòwán le} — I've finished.
\end{quote}

\begin{quote}
\emph{nǐ zuòwán le ma} — Have you finished?
\end{quote}

Its second character comes later; for now it is a word for the ear.

\begin{cousinweb}
\emph{zuòwán} is built from \zh{做} 'to do', the \zh{做} you write, and wán 'to finish'. \zh{做} has no reading in the Guangyun, so its history is not settled. wán is inherited from Middle Chinese \emph{hwan} (Guangyun, in Baxter's transcription), glossed 'whole; complete'; the regular sound change from \emph{hwan} would give huán, so why it is said wán is not settled.
\end{cousinweb}

\begin{grammarlens}[title={What you've built}]
\zh{做} is the doing; wán is the result: it is done. \zh{了} at the end makes it news: \emph{wǒ zuòwán le}.

Not finished yet takes \zh{没}: \emph{wǒ méi zuòwán}.
\end{grammarlens}

\subsection*{Your turn}
\begin{itemize}
\raggedright
  \item \emph{Say it:} \emph{zuòwán}
  \item \emph{Say it:} \emph{zuòwán}, once more
  \item \emph{Say it:} \emph{nǐ zuòwán le ma} — Have you finished?
  \item \emph{Say it:} \emph{wǒ zuòwán le} — I've finished.
  \item \emph{Recall:} say \zh{做饭}, then \emph{zuòwán}
\end{itemize}

\begin{tcolorbox}[breakable,colback=teal!4,colframe=teal!35!black,title={Before you move on}]
What does \emph{zuòwán} mean? (\textquotedblleft{}To finish doing\textquotedblright{}.) Say you haven't finished. (\emph{wǒ méi zuòwán}.)
\end{tcolorbox}
\section[Practice]{\zh{我看到他了}}

\label{lesson:ZH-C151-practice}



\begin{quote}
The chapter closes here, with two people checking on each other.

\begin{itemize}
\raggedright
  \item \emph{Recall:} say \emph{zuòwán}, to finish — \textbf{R1}, one lesson back
  \item \emph{Recall:} say \emph{kàndǒng}, to understand what you read — \textbf{R1}
  \item \emph{Recall:} say \zh{学会} — \emph{xuéhuì} — \textbf{R2}
\end{itemize}
\end{quote}

\subsection*{The exchange}
\begin{quote}
\textbf{A:} \zh{你看到他了吗}
\end{quote}

\begin{quote}
\emph{nǐ kàndào tā le ma} — Did you see him?
\end{quote}

\begin{quote}
\textbf{B:} \zh{看到了} \zh{他在那儿}
\end{quote}

\begin{quote}
\emph{kàndào le, tā zài nàr} — I did. He's over there.
\end{quote}

\begin{quote}
\textbf{A:} \zh{他说什么}
\end{quote}

\begin{quote}
\emph{tā shuō shénme} — What's he saying?
\end{quote}

\begin{quote}
\textbf{A:} \emph{nǐ tīngdǒng le ma} — Did you catch it?
\end{quote}

\begin{quote}
\textbf{B:} \emph{méi tīngdǒng, tā shuō de tài kuài le} — No. He talks too fast.
\end{quote}

\begin{quote}
\textbf{A:} \emph{nǐ zuòwán le ma} — Have you finished?
\end{quote}

\begin{quote}
\textbf{B:} \zh{没有} \zh{我在看书}
\end{quote}

\begin{quote}
\emph{méiyǒu, wǒ zài kàn shū} — No, I'm reading.
\end{quote}

\begin{quote}
\textbf{A:} \emph{zhè běn shū nǐ kàndǒng le ma} — Do you follow this book?
\end{quote}

\begin{quote}
\textbf{B:} \emph{kàndǒng le} — Yes.
\end{quote}

\begin{quote}
\textbf{A:} \zh{你学会开车了吗}
\end{quote}

\begin{quote}
\emph{nǐ xuéhuì kāichē le ma} — Have you learned to drive?
\end{quote}

\begin{quote}
\textbf{B:} \zh{学会了}
\end{quote}

\begin{quote}
\emph{xuéhuì le} — I have.
\end{quote}

Each yes repeats the verb and its result, with \zh{了}: \zh{看到了}, \emph{kàndǒng le}, \zh{学会了}. Each no puts \zh{没} in front: \emph{méi tīngdǒng}.

\begin{tcolorbox}[breakable,skin=enhanced,colback=orange!6,colframe=orange!45!black,boxrule=0.5pt,arc=1mm,left=6pt,right=6pt,top=4pt,bottom=4pt,fonttitle=\bfseries,title={Read this}]
\emph{Read it:} the exchange through once without stopping, then answer the question below
\end{tcolorbox}

\begin{grammarlens}[title={What you've built — five verbs that finish}]
A verb, then its result; \zh{没} when the result did not come:

\begin{itemize}
\raggedright
  \item \textbf{\zh{我看到他了}} \emph{wǒ kàndào tā le} — I saw him.
  \item \textbf{\zh{我学会开车了}} \emph{wǒ xuéhuì kāichē le} — I've learned to drive.
  \item \emph{wǒ méi tīngdǒng} — I didn't catch it.
  \item \emph{wǒ zuòwán le} — I've finished.
\end{itemize}

\textbf{Listen.} \emph{Nǐ kàndào wǒ de shǒujī le ma?}

\textbf{Speak.} Say one thing you have learned how to do, and one thing you have not finished.

\emph{Write it:} \zh{看到} — \zh{到} ends on the foot of \zh{在}, then the right of \zh{前}
\end{grammarlens}

\subsection*{Your turn}
\begin{itemize}
\raggedright
  \item \emph{Say it:} \zh{你看到他了吗} — \emph{nǐ kàndào tā le ma} — Did you see him?
  \item \emph{Say it:} \zh{看到了} — \emph{kàndào le} — I did.
  \item \emph{Say it:} the whole exchange, both parts
  \item \emph{Write it:} \zh{我看到他了}
\end{itemize}

\begin{tcolorbox}[breakable,colback=teal!4,colframe=teal!35!black,title={Before you move on}]
Say you saw him. (\textbf{\zh{我看到他了}}.) Say you didn't catch it. (\emph{wǒ méi tīngdǒng}.) Ask if someone has finished. (\emph{nǐ zuòwán le ma}.)
\end{tcolorbox}
